% !TeX program = xelatex
\documentclass[11pt,a4paper]{ctexart}
% Shared preamble for Agena architecture and process documents.
\usepackage[a4paper,margin=20mm]{geometry}
\usepackage{array}
\usepackage{booktabs}
\usepackage{caption}
\usepackage{enumitem}
\usepackage{fancyhdr}
\usepackage{hyperref}
\usepackage{longtable}
\usepackage{pdflscape}
\usepackage{tabularx}
\usepackage{tikz}
\usepackage[most]{tcolorbox}
\usepackage{xcolor}
\usepackage{listings}
\usepackage{url}

\usetikzlibrary{
  arrows.meta,
  backgrounds,
  calc,
  fit,
  positioning,
  shapes.geometric,
  shapes.misc,
  shapes.multipart
}

\definecolor{agenaBlue}{HTML}{1E5AA8}
\definecolor{agenaBlueSoft}{HTML}{EAF2FD}
\definecolor{agenaGreen}{HTML}{1D7A5E}
\definecolor{agenaGreenSoft}{HTML}{E7F6F1}
\definecolor{agenaOrange}{HTML}{B85C13}
\definecolor{agenaOrangeSoft}{HTML}{FCEFE4}
\definecolor{agenaPurple}{HTML}{5A4FCF}
\definecolor{agenaPurpleSoft}{HTML}{ECEAFB}
\definecolor{agenaInk}{HTML}{1D2433}
\definecolor{agenaMist}{HTML}{F7F9FC}
\definecolor{agenaSand}{HTML}{FFF9F0}
\definecolor{agenaGray}{HTML}{5D6678}
\definecolor{agenaBorder}{HTML}{C8D2E1}

\hypersetup{
  colorlinks=true,
  linkcolor=agenaBlue,
  urlcolor=agenaBlue,
  citecolor=agenaGreen,
  pdftitle={Agena Documentation},
  pdfauthor={OpenAI Codex}
}

\ctexset{
  section = {
    format = \Large\bfseries\color{agenaInk}
  },
  subsection = {
    format = \large\bfseries\color{agenaInk}
  },
  subsubsection = {
    format = \normalsize\bfseries\color{agenaInk}
  }
}

\setlength{\parskip}{0.55em}
\setlength{\parindent}{0pt}
\setlength{\headheight}{14.5pt}
\setlist[itemize]{leftmargin=2.2em}
\setlist[enumerate]{leftmargin=2.4em}
\captionsetup{hypcap=false}

\pagestyle{fancy}
\fancyhf{}
\fancyhead[L]{\textcolor{agenaGray}{Agena 文档}}
\fancyhead[R]{\textcolor{agenaGray}{\leftmark}}
\fancyfoot[C]{\textcolor{agenaGray}{\thepage}}
\renewcommand{\headrulewidth}{0.3pt}
\renewcommand{\footrulewidth}{0pt}

\newcommand{\docversion}{2026-04-06}
\newcommand{\pathinline}[1]{\texttt{\nolinkurl{#1}}}
\newcommand{\cmdinline}[1]{\texttt{#1}}
\newcommand{\term}[1]{\textbf{\textcolor{agenaInk}{#1}}}
\newcommand{\highlight}[1]{\textcolor{agenaBlue}{\textbf{#1}}}

\newtcolorbox{summarybox}[1][]{
  breakable,
  colback=agenaBlueSoft,
  colframe=agenaBlue,
  boxrule=0.8pt,
  arc=2pt,
  left=8pt,
  right=8pt,
  top=6pt,
  bottom=6pt,
  title=#1,
  fonttitle=\bfseries
}

\newtcolorbox{notebox}[1][]{
  breakable,
  colback=agenaSand,
  colframe=agenaOrange,
  boxrule=0.8pt,
  arc=2pt,
  left=8pt,
  right=8pt,
  top=6pt,
  bottom=6pt,
  title=#1,
  fonttitle=\bfseries
}

\newtcolorbox{techbox}[1][]{
  breakable,
  colback=agenaGreenSoft,
  colframe=agenaGreen,
  boxrule=0.8pt,
  arc=2pt,
  left=8pt,
  right=8pt,
  top=6pt,
  bottom=6pt,
  title=#1,
  fonttitle=\bfseries
}

\lstdefinestyle{agenaCode}{
  basicstyle=\ttfamily\small,
  backgroundcolor=\color{agenaMist},
  frame=single,
  rulecolor=\color{agenaBorder},
  columns=fullflexible,
  keepspaces=true,
  breaklines=true,
  showstringspaces=false
}

\tikzset{
  diagram/.style={
    font=\footnotesize,
    line join=round,
    line cap=round
  },
  module/.style={
    rectangle,
    rounded corners=4pt,
    draw=agenaInk!85,
    fill=white,
    thick,
    align=left,
    inner sep=8pt,
    text width=38mm,
    minimum height=12mm
  },
  modulewide/.style={
    module,
    text width=58mm
  },
  terminal/.style={
    module,
    draw=agenaBlue!80!black,
    fill=agenaBlueSoft
  },
  runtime/.style={
    module,
    draw=agenaGreen!80!black,
    fill=agenaGreenSoft
  },
  storage/.style={
    module,
    draw=agenaOrange!85!black,
    fill=agenaOrangeSoft
  },
  external/.style={
    module,
    draw=agenaPurple!80!black,
    fill=agenaPurpleSoft
  },
  note/.style={
    rectangle,
    rounded corners=4pt,
    draw=agenaBorder,
    fill=agenaMist,
    align=left,
    inner sep=6pt,
    text width=34mm,
    font=\scriptsize
  },
  decision/.style={
    diamond,
    draw=agenaInk!85,
    fill=agenaSand,
    aspect=2.2,
    align=center,
    inner sep=2pt,
    text width=24mm
  },
  group/.style={
    draw=agenaBorder,
    rounded corners=6pt,
    fill=agenaMist,
    inner sep=10pt
  },
  flow/.style={
    -{Latex[length=3mm,width=2mm]},
    thick,
    draw=agenaInk!85
  },
  signal/.style={
    flow,
    dashed,
    draw=agenaBlue!80!black
  },
  weakflow/.style={
    -{Latex[length=2.5mm,width=1.8mm]},
    semithick,
    draw=agenaGray
  },
  dbentity/.style={
    rectangle split,
    rectangle split parts=2,
    rectangle split part fill={agenaOrangeSoft,white},
    draw=agenaOrange!85!black,
    rounded corners=3pt,
    thick,
    align=left,
    text width=43mm,
    inner sep=6pt
  }
}

\newenvironment{diagramlandscape}
  {\clearpage\begin{landscape}\thispagestyle{plain}\begin{center}\captionsetup{width=0.94\linewidth}}
  {\end{center}\end{landscape}\clearpage}


\title{\textbf{Agena 架构文档}\\\large 面向小白的代码仓库架构说明}
\author{基于仓库源码整理}
\date{\docversion}

\begin{document}
\maketitle
\tableofcontents
\clearpage

\section{文档目标}

这份文档是给\highlight{第一次接触 Agena 仓库的人}准备的。

如果你以前没有看过 Rust Agent 后端，也没看过这个仓库，读完以后至少应该能回答下面五个问题：

\begin{enumerate}
  \item Agena 到底是一个什么系统，它解决什么问题？
  \item 这个仓库为什么要拆成 \pathinline{apps/}、\pathinline{crates/}、\pathinline{packages/}、\pathinline{ops/}？
  \item 真正的核心代码在哪，哪些只是“壳层”？
  \item 一次请求进入系统之后，会经过哪些核心模块？
  \item 如果我要继续读源码，应该从哪里开始看最不容易迷路？
\end{enumerate}

\begin{summarybox}[一句话先记住]
\term{Agena} 本质上是一个用 Rust 写的 \highlight{Agent 运行时后端}。它把\highlight{配置解析、Provider 适配、会话状态、工具调用、权限控制、插件扩展、数据库持久化}这些能力收在一个统一核心里，然后再由 CLI、HTTP API、Studio Server、Desktop 等外层应用去复用这套核心。
\end{summarybox}

\section{先建立一个正确心智模型}

很多新人第一次看这类仓库，最容易犯的错是把“页面、接口、桌面壳、CLI”都当成系统本体。对 Agena 来说，最好的理解方式恰好相反：

\begin{itemize}
  \item \highlight{真正的本体是 \pathinline{crates/agena}}。这里面有运行时、会话、消息、工具、Provider、权限、插件、数据库等核心能力。
  \item \pathinline{crates/agena-http-api} 是把核心能力包装成 HTTP API 的接口层。
  \item \pathinline{apps/agena-cli}、\pathinline{apps/agena-http-api-server}、\pathinline{apps/agena-studio-server} 是“可运行入口”，本质上是不同的启动壳。
  \item \pathinline{packages/agena-studio-web} 是前端页面，用来消费 Studio Server 暴露出来的 API。
  \item \pathinline{apps/agena-studio-desktop} 是 Tauri 桌面壳，用来把同一套 Web UI + 后端 sidecar 包成桌面应用。
\end{itemize}

\begin{notebox}[读仓库时最重要的顺序]
先看“核心”，再看“壳层”。更具体地说，建议顺序是：
\pathinline{crates/agena/src/lib.rs}
$\rightarrow$
\pathinline{runtime/}
$\rightarrow$
\pathinline{config/}
$\rightarrow$
\pathinline{provider/}
$\rightarrow$
\pathinline{session/}
$\rightarrow$
\pathinline{tool/}
$\rightarrow$
\pathinline{crates/agena-http-api}
$\rightarrow$
\pathinline{apps/agena-studio-server}
$\rightarrow$
\pathinline{packages/agena-studio-web}。
\end{notebox}

\section{仓库目录怎么读}

\begin{longtable}{p{0.22\textwidth}p{0.33\textwidth}p{0.35\textwidth}}
\toprule
\textbf{目录} & \textbf{作用} & \textbf{小白怎么读} \\
\midrule
\endhead
\pathinline{apps/} & 可运行应用入口。包括 CLI、HTTP API 服务、Studio Server、Desktop、TUI。 & 把它理解成“启动器”最合适。这里不是核心逻辑的主要沉淀地。 \\
\pathinline{crates/agena/} & 仓库真正的核心。包括 runtime、provider、session、tool、permission、plugin、db。 & 这是最值得深读的目录。绝大多数系统设计都在这里。 \\
\pathinline{crates/agena-http-api/} & 把核心能力变成资源化 HTTP API。 & 如果你想从前后端联调角度看系统，这里是桥梁层。 \\
\pathinline{packages/agena-studio-web/} & Studio 的 Vue 前端。负责聊天页、登录页、Runtime 页。 & 用来理解前端如何消费 API、如何用 SSE 做增量刷新。 \\
\pathinline{ops/} & 构建、打包、发布、安装相关脚本。 & 当你关心部署、桌面打包、服务安装时再看。 \\
\pathinline{docs/} & 项目文档和接口说明。 & 适合先扫一遍，但一定要和源码交叉验证。 \\
\bottomrule
\end{longtable}

\begin{diagramlandscape}
\begin{tikzpicture}[diagram, node distance=12mm and 14mm, scale=0.88, every node/.style={transform shape}]
  \node[terminal] (cli) {\textbf{agena-cli}\\读取配置并启动核心运行时\\目录：\pathinline{apps/agena-cli}};
  \node[terminal, below=of cli] (apiapp) {\textbf{agena-http-api-server}\\对外提供 REST + SSE\\目录：\pathinline{apps/agena-http-api-server}};
  \node[terminal, below=of apiapp] (studio) {\textbf{agena-studio-server}\\Studio API + UI 鉴权 + 静态资源\\目录：\pathinline{apps/agena-studio-server}};
  \node[terminal, below=of studio] (desktop) {\textbf{agena-studio-desktop}\\Tauri 桌面壳，管理 sidecar\\目录：\pathinline{apps/agena-studio-desktop}};

  \node[runtime, right=26mm of apiapp, text width=60mm, minimum height=16mm] (runtime) {\textbf{AgenaRuntime}\\运行时入口对象\\持有当前 \pathinline{RuntimeSnapshot}、后台任务控制、重载能力};

  \node[runtime, above right=14mm and 18mm of runtime] (config) {\textbf{Config 子系统}\\\pathinline{config/}\\默认值 + 文件 + mode + 环境变量 + CLI 覆盖};
  \node[runtime, right=22mm of runtime] (providers) {\textbf{ProviderRegistry}\\\pathinline{provider/}\\统一抽象 OpenAI / Anthropic / Gemini / Copilot / GitLab / Vertex 等};
  \node[runtime, below right=14mm and 18mm of runtime] (sessions) {\textbf{SessionManager}\\\pathinline{session/}\\会话状态、消息流、工具执行、阻塞恢复};
  \node[runtime, below left=14mm and 18mm of runtime] (plugins) {\textbf{PluginManager}\\\pathinline{plugin/}\\动态插件、hook、自定义 tool};
  \node[runtime, below=of sessions] (tools) {\textbf{ToolExecutor}\\\pathinline{tool/}\\bash / read / grep / apply\_patch / task / ask\_user};

  \node[storage, right=28mm of providers] (providerapi) {\textbf{外部模型服务}\\OpenAI / Anthropic / Gemini / GitLab / Vertex / Bedrock};
  \node[storage, below=of providerapi] (authstore) {\textbf{Auth Store}\\默认是文件存储\\保存 API key / OAuth token / 已知实例};
  \node[storage, below=of authstore] (db) {\textbf{SQLite + SeaORM}\\workspace / session / message / part / event / permission\_rule};
  \node[storage, below=of db] (workspacefs) {\textbf{本地工作区}\\Tool 读写文件，必要时走 \pathinline{procwarden} 沙箱};

  \begin{scope}[on background layer]
    \node[group, fit=(cli)(apiapp)(studio)(desktop), label={[text=agenaInk]north:可运行壳层}] {};
    \node[group, fit=(runtime)(config)(providers)(sessions)(plugins)(tools), label={[text=agenaInk]north:核心后端}] {};
    \node[group, fit=(providerapi)(authstore)(db)(workspacefs), label={[text=agenaInk]north:外部依赖与持久化}] {};
  \end{scope}

  \draw[flow] (cli.east) -- (runtime.west);
  \draw[flow] (apiapp.east) -- (runtime.west);
  \draw[flow] (studio.east) -- (runtime.west);
  \draw[flow] (desktop.east) -- node[above, font=\scriptsize] {启动 / 管理 sidecar} (studio.west);

  \draw[flow] (runtime) -- (config);
  \draw[flow] (runtime) -- (providers);
  \draw[flow] (runtime) -- (sessions);
  \draw[flow] (runtime) -- (plugins);
  \draw[flow] (sessions) -- (tools);
  \draw[flow] (providers) -- (providerapi);
  \draw[flow] (runtime) -- (authstore);
  \draw[flow] (sessions) -- (db);
  \draw[flow] (tools) -- (workspacefs);
  \draw[flow] (plugins) -- (tools);

  \node[note, right=10mm of desktop] (desktopnote) {桌面模式不是另一套后端。它本质上还是把 \pathinline{agena-studio-server} 当 sidecar 跑起来，只是外面包了一层 Tauri。};
\end{tikzpicture}
\captionof{figure}{Agena 仓库的总体系统架构。先分清“壳层”和“核心”，再去读局部代码就不容易迷路。}
\end{diagramlandscape}

\section{核心运行时层：AgenaRuntime 到底做了什么}

如果只允许你记一个类型名，那就是 \pathinline{AgenaRuntime}。

从源码看，\pathinline{AgenaRuntimeBuilder::build()} 会完成下面几件事：

\begin{enumerate}
  \item 决定工作区根目录。
  \item 连接数据库，并在需要时自动执行 migration。
  \item 通过 \pathinline{ConfigLoader} 解析配置。
  \item 构建第一代 \pathinline{RuntimeSnapshot}。
  \item 启动后台任务：
  \begin{itemize}
    \item \pathinline{reload::run}：轮询 watch path，看配置/插件是否变更。
    \item \pathinline{janitor::run}：定期清理 session cache。
  \end{itemize}
\end{enumerate}

这里有一个非常关键的设计点：\highlight{运行时是“快照化”的}。也就是说，系统不会把配置散落在各个单例里，而是把“当前有效配置 + 由配置派生出的服务对象”收拢成一个 \pathinline{RuntimeSnapshot}。这样做的好处是：

\begin{itemize}
  \item 手动 reload 或自动 reload 时，只需要重建新快照再整体替换。
  \item ProviderRegistry、PluginManager、SessionManager 等服务边界清晰。
  \item 外层应用只需要拿“当前快照”，不需要关心配置是怎么逐层覆盖出来的。
\end{itemize}

\begin{diagramlandscape}
\resizebox{0.94\linewidth}{!}{%
\begin{tikzpicture}[diagram, node distance=11mm and 14mm, scale=0.86, every node/.style={transform shape}]
  \node[runtime, text width=42mm] (loader) {\textbf{ConfigLoader}\\负责读取、合并、验证配置\\输出 \pathinline{ConfigResolution}};
  \node[runtime, right=of loader, text width=44mm] (resolution) {\textbf{ConfigResolution}\\最终配置 + 元信息\\active mode、config path、applied layers};
  \node[runtime, right=of resolution, text width=48mm] (snapshotbuild) {\textbf{RuntimeSnapshot::build}\\生成 ProviderRegistry、PluginManager、AuthStore、SessionManager};

  \node[runtime, below=18mm of snapshotbuild, text width=38mm] (store) {\textbf{RuntimeSnapshotStore}\\保存 current snapshot\\reload 时整体 swap};
  \node[runtime, right=of store, text width=42mm] (runtimebox) {\textbf{AgenaRuntime}\\对外暴露当前快照、reload、auth、session manager、workspace root};

  \node[storage, above left=13mm and 14mm of loader, text width=28mm] (defaults) {\textbf{内建默认值}\\代码里的默认配置};
  \node[storage, above=13mm of loader, text width=34mm] (file) {\textbf{config.toml}\\默认路径：\pathinline{~/.agena/config.toml}};
  \node[storage, above right=13mm and 14mm of loader, text width=28mm] (mode) {\textbf{mode 叠加}\\例如 \pathinline{dev} / \pathinline{prod}};
  \node[storage, below left=13mm and 14mm of loader, text width=30mm] (env) {\textbf{环境变量}\\例如 \pathinline{AGENA_MODE}};
  \node[storage, below right=13mm and 14mm of loader, text width=36mm] (cliovr) {\textbf{CLI 覆盖}\\例如 \pathinline{-c providers.openai.default_model=gpt-5}};

  \node[storage, above=13mm of snapshotbuild, text width=34mm] (workspace) {\textbf{workspace root}\\默认是当前工作目录};
  \node[storage, left=of store, text width=34mm] (dbconn) {\textbf{DatabaseConnection}\\SQLite 连接 + 自动 migration};
  \node[storage, below=13mm of store, text width=32mm] (watchpaths) {\textbf{watch paths}\\配置文件、插件目录等};
  \node[runtime, below=13mm of watchpaths, text width=36mm] (reloadtask) {\textbf{reload::run}\\轮询 watch paths\\发现变化后重建 snapshot};
  \node[runtime, below=13mm of runtimebox, text width=34mm] (janitor) {\textbf{janitor::run}\\定期修剪 session cache};

  \begin{scope}[on background layer]
    \node[group, fit=(defaults)(file)(mode)(env)(cliovr), label={[text=agenaInk]north:配置来源}] {};
    \node[group, fit=(loader)(resolution)(snapshotbuild)(store)(runtimebox)(reloadtask)(janitor), label={[text=agenaInk]north:运行时装配链}] {};
  \end{scope}

  \draw[flow] (defaults) -- (loader);
  \draw[flow] (file) -- (loader);
  \draw[flow] (mode) -- (loader);
  \draw[flow] (env) -- (loader);
  \draw[flow] (cliovr) -- (loader);
  \draw[flow] (loader) -- (resolution);
  \draw[flow] (resolution) -- (snapshotbuild);
  \draw[flow] (workspace) -- (snapshotbuild);
  \draw[flow] (dbconn.east) |- (snapshotbuild.south);
  \draw[flow] (snapshotbuild) -- (store);
  \draw[flow] (store) -- (runtimebox);
  \draw[flow] (watchpaths) -- (reloadtask);
  \draw[flow] (reloadtask) -- node[left, font=\scriptsize] {触发重建} (store.south);
  \draw[flow] (janitor) -- node[above, font=\scriptsize] {清理缓存} (store.east);
\end{tikzpicture}
}
\captionof{figure}{运行时装配链。Agena 不是把配置散落在各处，而是把“配置解析结果”和“由它派生的服务对象”聚合成一份可整体替换的 RuntimeSnapshot。}
\end{diagramlandscape}

\section{会话、消息、工具：真正最难也最重要的核心层}

对一个 Agent 后端来说，最复杂的永远不是“发个 HTTP 请求给模型”，而是：

\begin{itemize}
  \item 如何把一整个会话长期保存下来；
  \item 如何在工具调用、权限阻塞、用户补充输入之间来回切换；
  \item 如何把流式输出、结构化 part、数据库持久化、前端增量事件统一起来。
\end{itemize}

Agena 把这部分职责拆成了几个很清晰的角色：

\begin{itemize}
  \item \pathinline{SessionManager}
  \begin{itemize}
    \item 负责“调度循环”，核心方法是 \pathinline{run_until_stable()}。
    \item 它决定系统接下来是继续跑模型、执行工具、进入阻塞、还是已经 idle。
  \end{itemize}
  \item \pathinline{SessionStore}
  \begin{itemize}
    \item 负责数据库读写、缓存、ID 分配、事件落库。
  \end{itemize}
  \item \pathinline{SessionProcessor}
  \begin{itemize}
    \item 负责和 Provider 做真正的 completion / stream 通信。
    \item 它把 provider 流式事件投影成 message 和 part 的变更。
  \end{itemize}
  \item \pathinline{ToolExecutor}
  \begin{itemize}
    \item 负责真正执行 bash/read/apply\_patch/ask\_user 等工具。
  \end{itemize}
\end{itemize}

\begin{techbox}[这个拆分为什么好]
\pathinline{SessionManager} 只负责“决策下一步”，\pathinline{SessionProcessor} 只负责“跟模型交互并消费流”，\pathinline{SessionStore} 只负责“持久化和缓存”，\pathinline{ToolExecutor} 只负责“真正跑工具”。新人看代码时，如果先按这个职责边界去读，就不会在一个巨型函数里绕晕。
\end{techbox}

\begin{diagramlandscape}
\begin{tikzpicture}[diagram, node distance=12mm and 16mm, scale=0.95, every node/.style={transform shape}]
  \node[terminal] (http) {\textbf{HTTP API / Studio Server}\\接收 turns / continue / permission-replies / user-input-replies};
  \node[runtime, right=24mm of http, text width=58mm] (manager) {\textbf{SessionManager}\\\pathinline{submit_user_turn()}\\\pathinline{continue_session()}\\\pathinline{reply_permission()}\\\pathinline{reply_user_input()}\\核心循环：\pathinline{run_until_stable()}};

  \node[runtime, above right=14mm and 20mm of manager] (processor) {\textbf{SessionProcessor}\\Provider 流式调用\\MessageProjector\\prompt compaction / continuation};
  \node[runtime, below right=14mm and 20mm of manager] (executor) {\textbf{ToolExecutor}\\内建工具 + 插件工具\\沙箱、截断、tool catalog};
  \node[runtime, below=20mm of manager] (store) {\textbf{SessionStore}\\SeaORM CRUD\\session cache\\事件落库\\message / part ID 分配};
  \node[runtime, above=20mm of manager] (governor) {\textbf{ContextGovernor + prompt\_window}\\上下文预算、压缩、附件剥离、tool result pruning};

  \node[storage, right=22mm of processor] (providers) {\textbf{ProviderRegistry}\\根据 model ref 选择具体 Provider\\支持模型能力查询、流重放策略};
  \node[storage, below=of providers] (plugins) {\textbf{PluginManager}\\before\_tool / after\_tool / shell\_env / custom tool};
  \node[storage, below=of plugins] (policy) {\textbf{Permission 规则}\\默认策略 + 持久化 permission\_rules};
  \node[storage, below=of store] (sqlite) {\textbf{SQLite}\\session / message / part / detail / event / permission\_rule};
  \node[storage, right=22mm of executor] (workspace) {\textbf{本地工作区 + procwarden}\\真正发生文件读写、命令执行的位置};

  \begin{scope}[on background layer]
    \node[group, fit=(manager)(processor)(executor)(store)(governor), label={[text=agenaInk]north:会话核心}] {};
    \node[group, fit=(providers)(plugins)(policy)(workspace)(sqlite), label={[text=agenaInk]north:会话依赖}] {};
  \end{scope}

  \draw[flow] (http) -- (manager);
  \draw[flow] (manager) -- (processor);
  \draw[flow] (manager) -- (executor);
  \draw[flow] (manager) -- (store);
  \draw[flow] (manager) -- (governor);
  \draw[flow] (processor) -- (providers);
  \draw[flow] (executor) -- (plugins);
  \draw[flow] (executor) -- (policy);
  \draw[flow] (executor) -- (workspace);
  \draw[flow] (store) -- (sqlite);
  \draw[signal] (processor) to[bend left=10] node[above, font=\scriptsize] {assistant text / tool call / usage / finish} (store);
  \draw[signal] (store) to[bend left=10] node[below, font=\scriptsize] {loaded session / IDs / persisted state} (manager);
\end{tikzpicture}
\captionof{figure}{会话子系统的静态关系图。SessionManager 负责调度，Processor 负责模型流，Store 负责持久化，ToolExecutor 负责真正执行工具。}
\end{diagramlandscape}

\section{数据库与持久化模型}

从 entity 定义和 migration 可以看出，Agena 的数据库模型有两个重点：

\begin{itemize}
  \item \highlight{聊天历史不是一条大 JSON}，而是拆成 \pathinline{session}、\pathinline{message}、\pathinline{message_part}、\pathinline{message_part_detail} 四层。
  \item \highlight{会话状态不是只看 message 列表}，还要结合这些额外状态：
  \begin{itemize}
    \item \pathinline{runtime_state_json}
    \item \pathinline{session_events}
    \item \pathinline{permission_rules}
  \end{itemize}
\end{itemize}

这样设计的好处是：

\begin{itemize}
  \item message 列表可以做轻量分页；
  \item part 可以只返回 summary，不必每次都把 detail\_json 拉回前端；
  \item session event 可以单独做 SSE 增量流；
  \item session 的 \pathinline{version} 可以配合 HTTP \pathinline{If-Match} 做乐观并发控制。
\end{itemize}

\begin{diagramlandscape}
\begin{tikzpicture}[diagram, node distance=12mm and 14mm]
  \node[dbentity, text width=34mm] (workspace) {\textbf{agena\_workspaces}\nodepart{second}
    id\newline
    path\newline
    created\_at\_ms\newline
    updated\_at\_ms
  };

  \node[dbentity, text width=34mm, right=20mm of workspace] (session) {\textbf{agena\_sessions}\nodepart{second}
    id\newline
    parent\_id\newline
    workspace\_id\newline
    title\newline
    version\newline
    runtime\_state\_json\newline
    created\_at\_ms\newline
    updated\_at\_ms
  };

  \node[dbentity, text width=34mm, right=20mm of session] (message) {\textbf{agena\_messages}\nodepart{second}
    id\newline
    session\_id\newline
    role / status / source\newline
    parent\_message\_id\newline
    model\_provider\_id / model\_id\newline
    usage\_json / finish\newline
    tags\_json\newline
    created\_at\_ms / updated\_at\_ms
  };

  \node[dbentity, text width=34mm, right=20mm of message] (part) {{\scriptsize\textbf{agena\_message\_parts}}\nodepart{second}
    id\newline
    message\_id\newline
    part\_index\newline
    kind / status\newline
    summary\_text\newline
    has\_detail\newline
    call\_id / operation\_id\newline
    created\_at\_ms / updated\_at\_ms
  };

  \node[dbentity, text width=40mm, below=16mm of part] (detail) {{\scriptsize\textbf{agena\_message\_part}\\\textbf{details}}\nodepart{second}
    part\_id\newline
    detail\_json\newline
    updated\_at\_ms
  };

  \node[dbentity, text width=34mm, below=16mm of session] (event) {{\scriptsize\textbf{agena\_session\_events}}\nodepart{second}
    id\newline
    session\_id\newline
    seq\newline
    event\_type\newline
    payload\_json\newline
    causation\_id / correlation\_id\newline
    created\_at\_ms
  };

  \node[dbentity, text width=38mm, below=16mm of message] (rule) {{\scriptsize\textbf{agena\_permission}\\\textbf{rules}}\nodepart{second}
    id\newline
    action\_key\newline
    mode\newline
    created\_at\_ms\newline
    updated\_at\_ms
  };

  \draw[flow] (workspace) -- node[above] {\scriptsize 1:N} (session);
  \draw[flow] (session) -- node[above] {\scriptsize 1:N} (message);
  \draw[flow] (message) -- node[above] {\scriptsize 1:N} (part);
  \draw[flow] (part) -- node[right] {\scriptsize 1:0..1} (detail);
  \draw[flow] (session) -- node[left] {\scriptsize 1:N} (event);
  \draw[weakflow] (rule.north) -- ++(0,7mm) -| node[pos=0.3, above, font=\scriptsize] {按 action\_key 匹配工具权限} (session.south);

  \node[note, below=10mm of rule] (rulenote) {注意：\pathinline{permission_rules} 不直接挂在 session 上，而是按 action key 存持久化决策。这意味着“Allow always / Deny always”可以跨会话生效。};
\end{tikzpicture}
\captionof{figure}{Agena 的核心持久化模型。消息被拆成 message + part + detail，既便于懒加载，也便于保留结构化工具执行历史。}
\end{diagramlandscape}

\section{Studio 壳层与部署形态}

仓库里除了核心运行时，还有一条很明显的产品化路线，就是 \highlight{Agena Studio}。

Studio 不是重新实现了一套后端，而是在同一套核心上再加了：

\begin{itemize}
  \item \pathinline{agena-studio-server}：负责公开健康检查、UI 鉴权、CORS、静态资源服务、受保护 API。
  \item \pathinline{agena-studio-web}：Vue 前端页面，主要是 Chat / Runtime / Login。
  \item \pathinline{agena-studio-desktop}：Tauri 桌面壳，负责启动/停止 sidecar、读取桌面配置、托盘、更新等。
\end{itemize}

\begin{summarybox}[非常重要的一点]
Studio 的 Browser 模式和 Desktop 模式，\highlight{共享的是同一套后端资源语义}。不管你是浏览器打开还是桌面壳打开，本质上前端都还是在消费同一组 session / message / runtime / auth API。
\end{summarybox}

\begin{diagramlandscape}
\resizebox{0.90\linewidth}{!}{%
\begin{tikzpicture}[diagram, node distance=12mm and 16mm]
  \node[terminal, text width=48mm] (browser) {\textbf{浏览器模式}\\用户访问 Studio URL\\直接拿到 Web 静态资源};
  \node[terminal, below=of browser, text width=48mm] (desktopmode) {\textbf{桌面模式}\\Tauri 窗口 + 托盘 + 自动启动\\内部仍然加载同一套 Web UI};

  \node[terminal, right=22mm of browser, text width=52mm] (studioServer) {\textbf{agena-studio-server}\\公开路由：\pathinline{/health}、\pathinline{/auth/session}\\受保护路由：整套 \pathinline{/api/v1/*}\\可选静态资源服务：\pathinline{--ui-dir}};
  \node[terminal, right=22mm of desktopmode, text width=52mm] (desktopShell) {\textbf{agena-studio-desktop}\\读取桌面配置\\拉起 sidecar \pathinline{agena-studio}\\暴露 Tauri command 给前端};

  \node[storage, above right=12mm and 18mm of studioServer, text width=46mm] (webdist) {\textbf{Studio Web}\\目录：\pathinline{packages/agena-studio-web}\\构建产物：\pathinline{dist/}};
  \node[runtime, right=22mm of studioServer, text width=46mm] (uiauth) {\textbf{UI Auth}\\Studio 自己的会话密码层\\和 provider credential 是两回事};
  \node[runtime, below=14mm of uiauth, text width=56mm] (api) {\textbf{Agena HTTP API 资源语义}\\workspaces / sessions / messages / events / runtime / auth};
  \node[storage, right=22mm of api, text width=42mm] (runtimecore) {\textbf{AgenaRuntime}\\ProviderRegistry\\SessionManager\\PluginManager\\AuthStore};
  \node[storage, below=of runtimecore, text width=48mm] (persist) {\textbf{SQLite + 本地配置/日志}\\数据库、auth store、workspace};

  \begin{scope}[on background layer]
    \node[group, fit=(browser)(desktopmode), label={[text=agenaInk]north:用户入口}] {};
    \node[group, fit=(studioServer)(desktopShell), label={[text=agenaInk]north:Studio 壳层}] {};
    \node[group, fit=(webdist)(api)(uiauth)(runtimecore)(persist), label={[text=agenaInk]north:共享后端语义}] {};
  \end{scope}

  \draw[flow] (browser) -- (studioServer);
  \draw[flow] (desktopmode) -- (desktopShell);
  \draw[flow] (desktopShell) -- node[above, font=\scriptsize] {sidecar} (studioServer);
  \draw[flow] (webdist) -- (studioServer);
  \draw[flow] (studioServer) -- (uiauth);
  \draw[flow] (studioServer) |- (api);
  \draw[flow] (desktopShell) |- (api);
  \draw[flow] (api) -- (runtimecore);
  \draw[flow] (runtimecore) -- (persist);
\end{tikzpicture}
}
\captionof{figure}{Agena Studio 的两种运行形态。浏览器模式和桌面模式共享同一套后端资源语义，只是外层承载方式不同。}
\end{diagramlandscape}

\section{给小白的源码阅读顺序}

如果你打算继续往下读源码，我建议直接按这个顺序走：

\begin{enumerate}
  \item \pathinline{crates/agena/src/lib.rs}
  \item \pathinline{crates/agena/src/runtime/builder.rs}
  \item \pathinline{crates/agena/src/runtime/snapshot.rs}
  \item \pathinline{crates/agena/src/config/loader.rs}
  \item \pathinline{crates/agena/src/provider/registry.rs}
  \item \pathinline{crates/agena/src/session/manager.rs}
  \item \pathinline{crates/agena/src/session/processor.rs}
  \item \pathinline{crates/agena/src/session/store.rs}
  \item \pathinline{crates/agena/src/tool/mod.rs}
  \item \pathinline{crates/agena-http-api/src/lib.rs}
  \item \pathinline{apps/agena-studio-server/src/app.rs}
  \item \pathinline{packages/agena-studio-web/src/agena/pages/ChatPage.vue}
\end{enumerate}

\begin{notebox}[为什么是这个顺序]
因为这条路径是从“系统装配”往“会话执行”再往“接口和前端消费”推进的。它符合系统真正的依赖方向，也最容易在脑子里形成一条完整主线。
\end{notebox}

\section{本篇总结}

看完架构文档，你应该已经能把 Agena 分成三层：

\begin{itemize}
  \item \highlight{核心层}：\pathinline{crates/agena}
  \item \highlight{接口层}：\pathinline{crates/agena-http-api}
  \item \highlight{产品壳层}：\pathinline{apps/*} 和 \pathinline{packages/agena-studio-web}
\end{itemize}

下一篇《流程文档》会继续回答另一个关键问题：\highlight{一次用户请求，究竟是怎样一步一步在 Agena 里跑完的}。

\end{document}
